% ============================================================================
%  SECCIÓN 3.14: CONFIGURACIÓN DEL TEMA INICIAL
% ============================================================================

\section{Configuraci\'on del tema inicial}

La configuraci\'on del tema inicial define la identidad visual de la
aplicaci\'on, incluyendo el nombre visible, los iconos, el color de fondo y el
estilo de interfaz. En los proyectos de Expo, estos ajustes se declaran en el
archivo \texttt{app.json}, que centraliza la configuraci\'on de la aplicaci\'on
\citep{expo}.

El archivo \texttt{app.json} del proyecto se configur\'o con los siguientes
valores iniciales:

\begin{lstlisting}
{
  "expo": {
    "name": "lugares-interactivos",
    "slug": "lugares-interactivos",
    "version": "1.0.0",
    "orientation": "portrait",
    "icon": "./assets/icon.png",
    "userInterfaceStyle": "light",
    "android": {
      "adaptiveIcon": {
        "backgroundColor": "#E6F4FE"
      }
    }
  }
}
\end{lstlisting}

Los principales elementos del tema inicial se resumen en la tabla.

\begin{table}[H]
  \centering
  \caption{Elementos del tema inicial de la aplicaci\'on}
  \begin{tablaAPA}
    \begin{tabular}{@{}p{4.2cm}p{11.6cm}@{}}
      \toprule
      \textbf{Elemento} & \textbf{Descripci\'on} \\
      \midrule
      Nombre visible & \texttt{lugares-interactivos}, nombre provisional de la
      aplicaci\'on mostrado en el dispositivo. \\
      \addlinespace
      Orientaci\'on & Retrato (\texttt{portrait}), adecuada para el uso
      principal de la aplicaci\'on. \\
      \addlinespace
      Estilo de interfaz & Claro (\texttt{light}), como tema inicial del
      sistema \citep{expo}. \\
      \addlinespace
      Icono & \texttt{assets/icon.png}, imagen representativa de la
      aplicaci\'on. \\
      \addlinespace
      Color de fondo & \texttt{\#E6F4FE}, color claro que acompa\~na al icono
      adaptable en Android. \\
      \bottomrule
    \end{tabular}
  \end{tablaAPA}
\end{table}

\textbf{FUENTE:} Elaboraci\'on propia en base al repositorio del proyecto.

El tema inicial podr\'a evolucionar en avances posteriores, definiendo una
paleta de colores institucional, tipograf\'ias propias y un modo oscuro, en
funci\'on de los resultados de las pruebas de usabilidad previstas en el
proyecto.
